\chapter{O mundo quântico: fótons e níveis de energia}\label{ch:g12:quantum-world}

Uma lâmpada de calor infravermelha pode banhar a sua pele a noite inteira sem marcá-la;
dez minutos do sol rarefeito da montanha deixam uma queimadura. A lâmpada entrega muito
mais \omterm{def:g10:energy-conservation:energy}{energia} --- mas não é \omterm{def:g10:energy-conservation:energy}{energia} que marca a pele. A luz, mostra este capítulo, chega
em grãos, e só um grão que bate com \omterm{def:g10:inertia:force}{força} suficiente faz química: um grão \omterm{def:g10:light-spectra:wavelength}{ultravioleta}
rompe uma ligação que um bilhão de grãos \omterm{def:g10:light-spectra:wavelength}{infravermelhos} apenas aquecem. Siga os grãos e
um século se abre: metais vazando elétrons, \omterm{def:g11:fundamental-interactions:ladder}{átomos} cantando em \omterm{def:g10:light-spectra:emission}{linhas espectrais},
elétrons interferindo como \omterm{def:g10:signals-and-waves:wave}{ondas}.

\section{A luz chega em grãos: o fóton}

\begin{definition}[O fóton]\label{def:g12:quantum-world:photon}
A luz de \omterm{def:g12:oscillators-and-time:oscillator}{frequência} $\nu$ troca \omterm{def:g10:energy-conservation:energy}{energia} com a matéria apenas em grãos inteiros chamados
\emph{fótons}\index{fóton}, cada um de \omterm{def:g10:energy-conservation:energy}{energia}
\[
E = h\nu = \frac{hc}{\lambda}, \qquad h = \qty{6.63e-34}{J.s},
\]
onde $h$ é a \emph{constante de Planck}\index{constante de Planck} (Planck, 1900;
Einstein, 1905) e $\lambda$, o \omterm{def:g12:light-as-wave:lightwave}{comprimento de onda no vácuo}. A intensidade de um feixe
mede o \emph{número} de fótons por segundo; a \omterm{def:g10:energy-conservation:energy}{energia} de cada grão é fixada só pela
\omterm{def:g12:oscillators-and-time:oscillator}{frequência}.
\end{definition}

\begin{notation}[O elétron-volt]\label{not:g12:quantum-world:ev}
A física atômica usa o \emph{elétron-volt}\index{elétron-volt},
$\qty{1}{eV} = \qty{1.60e-19}{J}$: a \omterm{def:g10:energy-conservation:energy}{energia} que uma \omterm{def:g11:fundamental-interactions:elementary}{carga elementar} ganha ao atravessar
\qty{1}{V} (\cref{ch:g11:electric-gravitational-fields}). Prático:
$E \,(\unit{eV}) \approx 1240/\lambda\,(\unit{nm})$.
\end{notation}

\begin{omfigure}
\begin{tikzpicture}[font=\small, x=1.48cm]
  \draw[-Stealth] (7.55,0) -- (-0.55,0);
  \draw[-Stealth, omDef] (4.9,-0.85) -- (2.3,-0.85)
    node[midway, below] {a energia por fóton cresce};
  \foreach \x/\l/\e/\n in {0/{\qty{1}{pm}}/{\qty{1.2}{MeV}}/{$\gamma$},
      1.2/{\qty{0.1}{nm}}/{\qty{12}{keV}}/{X},
      2.4/{\qty{100}{nm}}/{\qty{12}{eV}}/{UV},
      3.6/{\qty{550}{nm}}/{\qty{2.3}{eV}}/{visível},
      4.8/{\qty{10}{\micro m}}/{\qty{0.12}{eV}}/{IV},
      6.0/{\qty{1}{cm}}/{\qty{0.12}{meV}}/{micro},
      7.2/{\qty{3}{m}}/{\qty{0.4}{\micro eV}}/{rádio}}{
    \draw (\x,0.07) -- (\x,-0.07);
    \node[above, omExo] at (\x,0.42) {\n};
    \node[above, font=\scriptsize] at (\x,0.10) {\l};
    \node[below, font=\scriptsize, omDef] at (\x,-0.10) {\e};}
\end{tikzpicture}

{\small Uma só escala para a família eletromagnética (\cref{ch:g12:light-as-wave}):
\omterm{def:g12:mechanical-waves:wavelength}{comprimento de onda} (em cima), \omterm{def:g10:energy-conservation:energy}{energia} por \omterm{def:g12:quantum-world:photon}{fóton} $hc/\lambda$ (em azul), dos grãos de
rádio de um \omterm{def:g11:circuits-and-power:voltage}{microelétron-volt} aos grãos gama de \unit{MeV} do
\cref{ch:g11:nucleus-radioactivity} --- e ainda assim um apontador de \qty{1.0}{mW}
despeja $\num{3e15}$ grãos por segundo: a luz do dia a dia escorre tão lisa quanto areia
derramada.}
\end{omfigure}

\section{O efeito fotoelétrico}

Ilumine uma placa metálica limpa no vácuo e elétrons saltam para fora --- o \emph{efeito
fotoelétrico}\index{efeito fotoelétrico}. As regras dele quebraram a física clássica.

\begin{proposition}[O que o experimento mostra]\label{prop:g12:quantum-world:facts}
Para cada metal existe uma \emph{frequência de corte}\index{frequência de corte} $\nu_0$:
\begin{enumerate}
  \item abaixo de $\nu_0$, elétron algum sai, por mais intensa que seja a luz;
  \item acima de $\nu_0$, a emissão começa sem \omterm{def:g12:mechanical-waves:delay}{atraso}, por mais fraca que ela seja;
  \item a intensidade eleva o \emph{número} de elétrons por segundo, e não a \omterm{def:g11:mechanical-energy:kinetic}{energia
        cinética} máxima deles; a \omterm{def:g12:oscillators-and-time:oscillator}{frequência} eleva $E_{k,\max}$.
\end{enumerate}
Nada disso cabe numa \omterm{def:g10:signals-and-waves:wave}{onda} contínua, com a qual qualquer cor deveria conseguir carregar
--- ao longo de meses, sob luz fraca (\cref{exo:g12:quantum-world:14}) --- e cuja
intensidade deveria fixar a \omterm{def:g10:energy-conservation:energy}{energia} dos elétrons: as três previsões falham.
\end{proposition}
\admitted

\begin{theorem}[Relação fotoelétrica de Einstein]\label{thm:g12:quantum-world:einstein}
\index{relação fotoelétrica}Um elétron precisa de pelo menos a \emph{função
trabalho}\index{função trabalho} $W_0$ do metal para escapar da superfície dele. Um
\omterm{def:g12:quantum-world:photon}{fóton} é absorvido inteiro por um elétron, de modo que
\[
h\nu = W_0 + E_{k,\max}, \qquad \text{e portanto } \nu_0 = \frac{W_0}{h}.
\]
\end{theorem}

\begin{proof}
Contabilidade de \omterm{def:g10:energy-conservation:energy}{energia}: entra $h\nu$, gasta-se $W_0$ na saída e o resto é cinético ---
no máximo $h\nu - W_0$, e menos para elétrons que partem de mais fundo. Abaixo de
$\nu_0$ um grão não consegue pagar o pedágio de saída, seja qual for o número deles.
\end{proof}

\begin{definition}[Potencial de corte]\label{def:g12:quantum-world:stopping}
Para medir $E_{k,\max}$, uma tensão reversa freia os elétrons
(\cref{ch:g11:electric-gravitational-fields}); o \emph{potencial de
corte}\index{potencial de corte} $U_s$ é a menor tensão que anula a corrente: os
elétrons mais rápidos por pouco não conseguem atravessar, de modo que
$E_{k,\max} = e\,U_s$.
\end{definition}

\begin{omfigure}
\begin{tikzpicture}
\begin{axis}[omaxis, width=8.8cm, height=5.4cm,
    xlabel={$\nu$ ($10^{14}\,\unit{Hz}$)}, ylabel={$E_{k,\max}$ (\unit{eV})},
    xmin=0, xmax=13.4, ymin=-3.1, ymax=3.2,
    xtick={2,4,6,8,10,12}, ytick={-2,-1,1,2,3}]
  \addplot[omExo, dashed, domain=0:5.55, samples=2] {0.414*x-2.3};
  \addplot[omThm, thick, domain=5.55:12.4, samples=2] {0.414*x-2.3};
  \addplot[only marks, mark=*, mark size=1.5pt, omDef] coordinates
    {(7,0.60) (8,1.01) (9,1.43) (10,1.84) (11,2.26)};
  \node[below, font=\small] at (axis cs:5.55,-0.08) {$\nu_0$};
  \node[right, font=\small, omExo] at (axis cs:0.15,-2.55) {$-W_0$};
  \node[font=\small, omThm, rotate=22] at (axis cs:10.3,1.35) {inclinação $= h$};
\end{axis}
\end{tikzpicture}

{\small $E_{k,\max} = eU_s$ medido contra a \omterm{def:g12:oscillators-and-time:oscillator}{frequência} (sódio): uma reta de inclinação
$h$, cortando o eixo em $\nu_0$; extrapolada (tracejado), ela revela $-W_0$ em
$\nu = 0$.}
\end{omfigure}

\begin{method}[Medir a constante de Planck]\label{met:g12:quantum-world:measure-h}
\begin{enumerate}
  \item Ilumine uma célula fotoelétrica através de \omterm{def:g11:color-light-sources:subtractive}{filtros} de \omterm{def:g12:mechanical-waves:wavelength}{comprimentos de onda}
        conhecidos; calcule cada $\nu = c/\lambda$; registre cada \omterm{def:g12:quantum-world:stopping}{potencial de corte}.
  \item Trace $E_{k,\max} = eU_s$ contra $\nu$: Einstein promete a reta $h\nu - W_0$.
  \item Leia a inclinação: ela \emph{é} $h$; os interceptos dão $\nu_0$ e $-W_0$.
\end{enumerate}
\end{method}

\section{Níveis de energia: por que os átomos emitem linhas}

\begin{definition}[Níveis de energia]\label{def:g12:quantum-world:levels}
A \omterm{def:g10:energy-conservation:energy}{energia} de um elétron ligado num \omterm{def:g11:fundamental-interactions:ladder}{átomo} é \emph{quantizada}\index{quantização}:
restrita a uma escada discreta de \emph{níveis de energia}\index{nível de energia}
$E_1 < E_2 < \dots < 0$, contados negativos a partir do elétron livre. O degrau mais
baixo é o \emph{estado fundamental}\index{estado fundamental}, e os outros são
\emph{estados excitados}\index{estado excitado}; chegar a $0$ é a
\emph{ionização}\index{ionização}.
\end{definition}

\begin{proposition}[A escada do hidrogênio]\label{prop:g12:quantum-world:hydrogen}
Os \omterm{def:g12:quantum-world:levels}{níveis de energia} do \omterm{def:g11:fundamental-interactions:ladder}{átomo} de hidrogênio são
\[
E_n = -\frac{\qty{13.6}{eV}}{n^2}, \qquad n = 1, 2, 3, \dots
\]
de modo que $E_1 = \qty{-13.6}{eV}$, $E_2 = \qty{-3.40}{eV}$,
$E_3 = \qty{-1.51}{eV}$, $E_4 = \qty{-0.85}{eV}$, amontoando-se rumo a $0$.
\end{proposition}
\admitted

\begin{remark}[De onde vem a escada]\label{rem:g12:quantum-world:ladder-origin}
Por que esses números é mecânica quântica, guardada para os volumes universitários ---
assim como a fórmula de de Broglie: a \omterm{def:g12:quantum-world:debroglie}{onda de matéria} do elétron
(\cref{def:g12:quantum-world:debroglie}, adiante) precisa se fechar sobre si mesma em
torno do \omterm{def:g11:fundamental-interactions:ladder}{núcleo}, e só certos \omterm{def:g12:mechanical-waves:wavelength}{comprimentos de onda} cabem --- como uma corda de violão, um
\omterm{def:g11:fundamental-interactions:ladder}{átomo} tem um conjunto discreto de notas.
\end{remark}

\begin{theorem}[O fóton de uma transição]\label{thm:g12:quantum-world:transition}
Um \omterm{def:g11:fundamental-interactions:ladder}{átomo} muda de nível apenas emitindo ou absorvendo um \omterm{def:g12:quantum-world:photon}{fóton} que carrega exatamente a
diferença: cair do nível $p$ para um nível mais baixo $n$ emite $h\nu = E_p - E_n$;
subir de $n$ até $p$ absorve essa mesma \omterm{def:g10:energy-conservation:energy}{energia}.
\end{theorem}

\begin{proof}
\omterm{thm:g10:energy-conservation:conservation}{Conservação da energia}, um grão de cada vez: o \omterm{def:g12:quantum-world:photon}{fóton} é a única moeda que a luz aceita, e
os livros do \omterm{def:g11:fundamental-interactions:ladder}{átomo} precisam fechar.
\end{proof}

\begin{remark}[Espectros de linhas explicados]\label{rem:g12:quantum-world:lines}
Aqui está a resposta prometida pelo \cref{ch:g10:light-spectra}: um gás quente emite
apenas as \omterm{def:g12:oscillators-and-time:oscillator}{frequências} $(E_p - E_n)/h$ da escada dele --- um código de barras de linhas
brilhantes --- e um gás frio rouba exatamente essas \omterm{def:g12:oscillators-and-time:oscillator}{frequências} da \omterm{def:g10:light-spectra:white}{luz branca},
imprimindo linhas escuras de absorção que casam com elas. Cada elemento tem a escada
dele, e portanto o código de barras dele --- é assim que lemos a luz das estrelas.
\end{remark}

\begin{example}[As linhas visíveis do hidrogênio]\label{ex:g12:quantum-world:balmer}
As transições que terminam em $n = 2$ caem no visível. Com
$E \,(\unit{eV}) = 1240/\lambda\,(\unit{nm})$: $3 \to 2$: \qty{1.89}{eV}, \qty{656}{nm}
(vermelho); $4 \to 2$: \qty{2.55}{eV}, \qty{486}{nm} (verde-azulado); $5 \to 2$:
\qty{2.86}{eV}, \qty{434}{nm} (violeta). As transições para $n = 1$ passam de
\qty{10}{eV}: todas \omterm{def:g10:light-spectra:wavelength}{ultravioleta}.
\end{example}

\begin{omfigure}
\begin{tikzpicture}[font=\small, y=1cm]
  \foreach \y/\n/\e in {0/1/-13.6, 4.50/2/-3.40, 5.33/3/-1.51}{
    \draw[thick] (0,\y) -- (4.2,\y);
    \node[left] at (0,\y) {$n=\n$};
    \node[right, font=\scriptsize] at (4.2,\y) {\qty{\e}{eV}};}
  \draw[thick] (0,5.63) -- (4.2,5.63) node[left,at start,yshift=-2.5pt] {$n=4$};
  \draw[thick] (0,5.76) -- (4.2,5.76) node[left,at start,yshift=3.5pt] {$n=5$};
  \draw[dashed] (0,6.0) -- (4.2,6.0) node[left, at start] {$\infty$}
    node[right, font=\scriptsize] {\qty{0}{eV} (ionizado)};
  \foreach \x/\ys/\col/\l in {1.0/5.33/omThm/656, 1.9/5.63/omDef/486,
      2.8/5.76/violet/434} \draw[-Stealth, \col, thick] (\x,\ys) --
    (\x,4.50) node[midway, left, font=\scriptsize] {\qty{\l}{nm}};
  \draw[-Stealth, omExo, thick] (3.7,4.50) -- (3.7,0)
    node[midway, right, font=\scriptsize] {\qty{122}{nm} (UV)};
\end{tikzpicture}

{\small A escada do hidrogênio, em escala: as linhas visíveis terminam todas em $n = 2$;
a queda até o \omterm{def:g12:quantum-world:levels}{estado fundamental} é \omterm{def:g10:light-spectra:wavelength}{ultravioleta} profundo.}
\end{omfigure}

\begin{remark}[Lasers, num fôlego só]\label{rem:g12:quantum-world:laser}
Um \omterm{def:g12:quantum-world:photon}{fóton} de exatamente $E_p - E_n$ que passe por um \omterm{def:g11:fundamental-interactions:ladder}{átomo} \emph{excitado} pode
provocá-lo a emitir um \omterm{def:g12:quantum-world:photon}{fóton} idêntico --- mesma \omterm{def:g10:energy-conservation:energy}{energia}, mesma direção, em compasso.
Mantenha mais \omterm{def:g11:fundamental-interactions:ladder}{átomos} em cima do que embaixo e um \omterm{def:g12:quantum-world:photon}{fóton} vira uma avalanche de clones:
\emph{emissão estimulada}\index{emissão estimulada} --- e um \omterm{def:g11:color-light-sources:lamps}{laser} é pouco mais do que
isso.
\end{remark}

\section{Ondas de matéria}

\begin{proposition}[Os elétrons interferem]\label{prop:g12:quantum-world:interference}
Mande elétrons um a um por uma fenda dupla (\cref{ch:g12:light-as-wave}): cada um chega
como um ponto nítido --- uma partícula. Mas, à medida que milhares se acumulam, os
pontos desenham \omterm{def:g12:light-as-wave:fringes}{franjas de interferência}: cada elétron, sozinho, passa como \omterm{def:g10:signals-and-waves:wave}{onda} pelas
\emph{duas} fendas. Cristais também difratam feixes de elétrons (Davisson e Germer,
1927).
\end{proposition}
\admitted

\begin{omfigure}
\begin{tikzpicture}[font=\small]
  \foreach \o in {0,4.0,8.0} \draw[omExo] (\o,0) rectangle +(3.4,2.2);
  \foreach \x/\y in {0.62/1.71, 1.68/0.42, 2.31/1.95, 1.05/1.22,
      2.98/0.66, 1.74/1.58, 0.38/0.35, 2.40/0.98, 1.62/1.02, 0.97/1.86,
      2.95/1.60, 1.71/0.20, 1.12/0.61, 2.33/0.31}
    \fill (\x,\y) circle (0.8pt);
  \foreach \c/\f in {0.40/610, 1.05/410, 1.70/530, 2.35/350, 3.00/470}
    \foreach \y in {0.25,0.62,...,2.00}
      \fill ({4.0+\c+0.14*sin(\f*\y)},\y) circle (0.8pt);
  \foreach \x/\y in {4.72/1.42, 5.38/0.52, 6.06/1.12, 6.66/1.72,
      5.44/1.88, 7.28/0.78} \fill (\x,\y) circle (0.8pt);
  \foreach \c/\f in {0.40/610, 1.05/410, 1.70/530, 2.35/350, 3.00/470}
    \foreach \y in {0.12,0.27,...,2.10}
      \fill ({8.0+\c+0.11*sin(\f*\y)},\y) circle (0.8pt);
  \foreach \c/\f in {0.40/260, 1.05/340, 1.70/205, 2.35/290, 3.00/385}
    \foreach \y in {0.20,0.50,...,2.05}
      \fill ({8.0+\c+0.06*sin(\f*\y+40)},\y) circle (0.8pt);
  \node[below] at (1.7,-0.05) {$\sim 30$ elétrons};
  \node[below] at (5.7,-0.05) {$\sim 200$}; \node[below] at (9.7,-0.05) {$\sim 2000$};
\end{tikzpicture}

{\small Um elétron, um ponto; milhares, franjas. Onde o próximo ponto cai é aleatório ---
como o decaimento de um \omterm{def:g11:fundamental-interactions:ladder}{núcleo} (\cref{ch:g11:nucleus-radioactivity}) --- mas a \omterm{def:g10:signals-and-waves:wave}{onda} fixa
a \emph{probabilidade}: franja clara, provável; escura, quase nunca. O acaso individual
governado por uma \omterm{def:g10:signals-and-waves:wave}{onda} exata: o coração da física quântica.}
\end{omfigure}

\begin{definition}[Comprimento de onda de de Broglie]\label{def:g12:quantum-world:debroglie}
Toda partícula de \omterm{def:g12:newtons-laws:momentum}{quantidade de movimento} $p = mv$ viaja como uma \emph{onda de
matéria}\index{onda de matéria} de \emph{comprimento de onda de de
Broglie}\index{comprimento de onda de de Broglie}
\[
\lambda = \frac{h}{p} = \frac{h}{mv}.
\]
\end{definition}
\admitted

\begin{example}[Por que você nunca difrata]\label{ex:g12:quantum-world:never}
Um elétron acelerado por \qty{100}{V} tem $E_k = \qty{1.60e-17}{J}$ e
$p = \sqrt{2m_e E_k} = \qty{5.4e-24}{kg.m/s}$, de modo que
$\lambda \approx \qty{0.12}{nm}$ --- o espaçamento dos \omterm{def:g11:fundamental-interactions:ladder}{átomos} num cristal, que por isso
funciona como a fenda dupla dele. Uma bola de tênis de \qty{58}{g} a \qty{25}{m/s}:
$\lambda = 6.63\times10^{-34}/1.45 \approx \qty{4.6e-34}{m}$, umas $10^{19}$ vezes menor
do que um \omterm{def:g11:fundamental-interactions:ladder}{núcleo} --- fenda alguma jamais a mostrará em franjas. O $h$ é tão pequeno que
a granulosidade quântica só vem à tona na escala do \omterm{def:g11:fundamental-interactions:ladder}{átomo}.
\end{example}

\begin{example}[O microscópio eletrônico]\label{ex:g12:quantum-world:microscope}
A \omterm{def:g12:light-as-wave:diffraction}{difração} proíbe resolver detalhes muito menores do que o \omterm{def:g12:mechanical-waves:wavelength}{comprimento de onda} usado
(\cref{ch:g12:light-as-wave}): a luz visível para perto de \qty{0.5}{\micro m}. Elétrons
a \qty{10}{kV} têm $\lambda \approx \qty{12}{pm}$ --- dezenas de milhares de vezes menor:
os microscópios eletrônicos fotografam vírus, a maquinaria da célula e até \omterm{def:g11:fundamental-interactions:ladder}{átomos}
isolados. A \omterm{def:g12:quantum-world:debroglie}{onda de matéria} entregou à física um olho novo.
\end{example}

\section{Exercícios}

\begin{exercise}[$\star$]\label{exo:g12:quantum-world:1}
Um apontador \omterm{def:g11:color-light-sources:lamps}{laser} verde emite \qty{1.0}{mW} em \qty{530}{nm}. Calcule a \omterm{def:g10:energy-conservation:energy}{energia} de um
\omterm{def:g12:quantum-world:photon}{fóton} em \omterm{def:g11:work-of-force:work}{joules} e em \omterm{def:g12:nuclear-energy:units}{elétron-volts}, e depois os \omterm{def:g12:quantum-world:photon}{fótons} por segundo.
\end{exercise}

\begin{exercise}[$\star$]\label{exo:g12:quantum-world:2}
Uma estação de FM transmite a \qty{100}{MHz}. Calcule a \omterm{def:g10:energy-conservation:energy}{energia} do \omterm{def:g12:quantum-world:photon}{fóton} dela em
\unit{eV} e compare com a de um \omterm{def:g12:quantum-world:photon}{fóton} visível: por que a granulosidade do rádio é
indetectável?
\end{exercise}

\begin{exercise}[$\star$]\label{exo:g12:quantum-world:3}
O césio tem $W_0 = \qty{1.9}{eV}$. Calcule a \omterm{def:g12:oscillators-and-time:oscillator}{frequência} e o \omterm{def:g12:mechanical-waves:wavelength}{comprimento de onda} de corte
dele. Um \omterm{def:g11:color-light-sources:lamps}{laser} vermelho de \qty{633}{nm} arranca elétrons do césio --- e, se arranca, com
que \omterm{def:g11:mechanical-energy:kinetic}{energia cinética} máxima?
\end{exercise}

\begin{exercise}[$\star$]\label{exo:g12:quantum-world:4}
Usando $E_n = -13.6/n^2$ \unit{eV}, calcule a \omterm{def:g10:energy-conservation:energy}{energia} e o \omterm{def:g12:mechanical-waves:wavelength}{comprimento de onda} do \omterm{def:g12:quantum-world:photon}{fóton}
emitido na transição $3 \to 2$ do hidrogênio. Qual é a cor dele?
\end{exercise}

\begin{exercise}[$\star$]\label{exo:g12:quantum-world:5}
No gráfico de $E_{k,\max}$ contra $\nu$ do curso, leia a \omterm{prop:g12:quantum-world:facts}{frequência de corte} e a \omterm{thm:g12:quantum-world:einstein}{função
trabalho} do sódio (em \unit{eV}); que constante é a inclinação?
\end{exercise}

\begin{exercise}[$\star\star$]\label{exo:g12:quantum-world:6}
Luz de \omterm{def:g12:mechanical-waves:wavelength}{comprimento de onda} \qty{400}{nm} incide sobre potássio
($W_0 = \qty{2.3}{eV}$). Calcule a \omterm{def:g10:energy-conservation:energy}{energia} do \omterm{def:g12:quantum-world:photon}{fóton}, a \omterm{def:g11:mechanical-energy:kinetic}{energia cinética} máxima dos
elétrons (em \unit{eV} e em \omterm{def:g11:work-of-force:work}{joules}) e o \omterm{def:g12:quantum-world:stopping}{potencial de corte}.
\end{exercise}

\begin{exercise}[$\star\star$]\label{exo:g12:quantum-world:7}
Uma célula fotoelétrica é iluminada acima do corte. O que acontece com (a) a corrente e
(b) o \omterm{def:g12:quantum-world:stopping}{potencial de corte} quando a intensidade é dobrada com a \omterm{def:g12:oscillators-and-time:oscillator}{frequência} fixa? E quando
a \omterm{def:g12:oscillators-and-time:oscillator}{frequência} é elevada? Que fato contradiz a imagem ondulatória, e por quê?
\end{exercise}

\begin{exercise}[$\star\star$]\label{exo:g12:quantum-world:8}
Que \omterm{def:g10:energy-conservation:energy}{energia} mínima de \omterm{def:g12:quantum-world:photon}{fóton} ioniza um \omterm{def:g11:fundamental-interactions:ladder}{átomo} de hidrogênio a partir do \omterm{def:g12:quantum-world:levels}{estado
fundamental}? E a partir de $n = 2$? Calcule os dois \omterm{def:g12:mechanical-waves:wavelength}{comprimentos de onda} de corte e
nomeie os domínios espectrais deles.
\end{exercise}

\begin{exercise}[$\star\star$]\label{exo:g12:quantum-world:9}
\omterm{def:g10:light-spectra:white}{Luz branca} atravessa um frasco de hidrogênio frio (todos os \omterm{def:g11:fundamental-interactions:ladder}{átomos} no \omterm{def:g12:quantum-world:levels}{estado
fundamental}). Explique por que o \omterm{def:g12:sound-acoustics:timbre}{espectro} transmitido \emph{não} mostra linha escura
alguma no visível, e ainda assim as linhas de absorção de Balmer aparecem nos \omterm{def:g12:sound-acoustics:timbre}{espectros}
de estrelas quentes.
\end{exercise}

\begin{exercise}[$\star\star$]\label{exo:g12:quantum-world:10}
Calcule o \omterm{def:g12:quantum-world:debroglie}{comprimento de onda de de Broglie} (a) de um elétron acelerado por
\qty{100}{V}; (b) de uma mosca de \qty{0.10}{g} a \qty{1.0}{m/s}. Compare cada um com o
tamanho de um \omterm{def:g11:fundamental-interactions:ladder}{átomo} ($\sim \qty{e-10}{m}$); conclua.
\end{exercise}

\begin{exercise}[$\star\star$]\label{exo:g12:quantum-world:11}
Um olho adaptado ao escuro responde a cerca de $90$ \omterm{def:g12:quantum-world:photon}{fótons} de luz de \qty{505}{nm}.
Calcule essa \omterm{def:g10:energy-conservation:energy}{energia}, compare-a com a \omterm{def:g11:mechanical-energy:kinetic}{energia cinética} de um mosquito de \qty{2.5}{mg}
voando a \qty{0.50}{m/s} e interprete a razão.
\end{exercise}

\begin{exercise}[$\star\star\star$]\label{exo:g12:quantum-world:12}
Uma célula fotoelétrica dá $U_s = \qty{1.19}{V}$ em \qty{405}{nm} e
$U_s = \qty{0.40}{V}$ em \qty{546}{nm}. A partir desses dois pontos, meça $h$ e depois a
\omterm{thm:g12:quantum-world:einstein}{função trabalho} em \unit{eV}. Que metal deste capítulo é o cátodo?
\end{exercise}

\begin{exercise}[$\star\star\star$]\label{exo:g12:quantum-world:13}
Uma lâmpada de hidrogênio mostra linhas em \qty{486}{nm} e \qty{434}{nm}. Converta cada
uma numa \omterm{def:g10:energy-conservation:energy}{energia} de \omterm{def:g12:quantum-world:photon}{fóton} e identifique as duas transições casando diferenças dos níveis
$E_n = -13.6/n^2$ \unit{eV}.
\end{exercise}

\begin{exercise}[$\star\star\star$]\label{exo:g12:quantum-world:14}
Estimativa clássica do \omterm{def:g12:mechanical-waves:delay}{atraso}: luz fraca, de intensidade \qty{1.0e-6}{W/m^2}, incide
sobre potássio ($W_0 = \qty{2.3}{eV}$). Se um \omterm{def:g11:fundamental-interactions:ladder}{átomo} coletasse apenas a \omterm{def:g10:energy-conservation:energy}{energia} que
atravessa a seção reta dele (raio \qty{1.0e-10}{m}), quanto tempo levaria para juntar
$W_0$? Compare com o \omterm{def:g12:mechanical-waves:delay}{atraso} observado (menos de \qty{e-9}{s}).
\end{exercise}

\begin{exercise}[$\star\star\star$]\label{exo:g12:quantum-world:15}
Num microscópio eletrônico os elétrons são acelerados por \qty{5.0}{kV}. Calcule o
\omterm{def:g12:quantum-world:debroglie}{comprimento de onda de de Broglie} deles, compare-o com o da luz verde (\qty{550}{nm}) e
explique, pela \omterm{def:g12:light-as-wave:diffraction}{difração} (\cref{ch:g12:light-as-wave}), por que isso compra resolução.
\end{exercise}

\section{Problema: o laboratório de física do painel solar}

\begin{problem}[{Problema de fim de semana --- o laboratório de física do painel solar:
os $10^{21}$ grãos que um telhado bebe a cada segundo, uma célula fotoelétrica que mede
a \omterm{def:g12:quantum-world:photon}{constante de Planck}, uma lâmpada de hidrogênio lida linha a linha --- um fim de semana,
um $h$}]
\label{pb:g12:quantum-world:1}
A sua turma instrumenta o painel solar do telhado da escola por um fim de semana. Dados:
irradiância solar ao meio-dia \qty{1.0e3}{W/m^2}; \omterm{def:g12:mechanical-waves:wavelength}{comprimento de onda} solar médio
\qty{550}{nm}; área do painel \qty{1.0}{m^2}, \omterm{def:g10:energy-conservation:efficiency}{rendimento} $20\%$; o dia entrega o
equivalente a \qty{5.0}{h} de sol pleno.

\medskip
\noindent\textbf{Parte I --- Contabilidade de \omterm{def:g12:quantum-world:photon}{fótons}.}
\begin{enumerate}
  \item Calcule a \omterm{def:g10:energy-conservation:energy}{energia} de um \omterm{def:g12:quantum-world:photon}{fóton} solar médio, em \omterm{def:g11:work-of-force:work}{joules} e em \omterm{def:g12:nuclear-energy:units}{elétron-volts}.
  \item Deduza o número de \omterm{def:g12:quantum-world:photon}{fótons} que atingem o painel por segundo ao meio-dia.
  \item Calcule a \omterm{prop:g11:circuits-and-power:power}{potência elétrica} do painel ao meio-dia e depois a \omterm{def:g10:energy-conservation:energy}{energia} produzida ao
        longo do dia, em \unit{kW.h}.
  \item A corrente do painel parece perfeitamente lisa. Usando a questão 2, explique por
        que a chegada grão a grão é invisível.
  \item A luz do Sol carrega mais \omterm{def:g10:energy-conservation:energy}{energia} infravermelha do que \omterm{def:g10:light-spectra:wavelength}{ultravioleta} e, ainda
        assim, só a \omterm{def:g10:light-spectra:wavelength}{ultravioleta} queima a pele (ligações de $\qtyrange{3}{4}{eV}$).
        Explique em termos de grãos.
\end{enumerate}

\medskip
\noindent\textbf{Parte II --- Medir $h$ com a célula fotoelétrica do kit.}
A célula fotoelétrica do kit tem um cátodo desconhecido. Os \omterm{def:g11:color-light-sources:subtractive}{filtros} dão:
\begin{center}\small
\begin{tabular}{lcccc}
$\lambda$ (\unit{nm}) & 365 & 405 & 436 & 546\\
$U_s$ (\unit{V}) & 1.50 & 1.16 & 0.94 & 0.37
\end{tabular}
\end{center}
\begin{enumerate}[resume]
  \item O que o \omterm{def:g12:quantum-world:stopping}{potencial de corte} mede? Relacione $U_s$, $\nu$, $W_0$ e $h$.
  \item Converta os quatro \omterm{def:g12:mechanical-waves:wavelength}{comprimentos de onda} em \omterm{def:g12:oscillators-and-time:oscillator}{frequências} (em
        $10^{14}\,\unit{Hz}$).
  \item Por que traçar $eU_s$ contra $\nu$ deve dar uma reta? Verifique o alinhamento
        dos pontos.
  \item A partir dos pontos extremos, calcule a inclinação: o seu $h$ medido.
  \item Deduza a \omterm{thm:g12:quantum-world:einstein}{função trabalho} do cátodo (em \unit{eV}), a \omterm{prop:g12:quantum-world:facts}{frequência de corte} e o
        \omterm{def:g12:mechanical-waves:wavelength}{comprimento de onda} de corte.
  \item Que parte da luz do Sol conseguiria arrancar elétrons desse cátodo --- e que
        metade do \omterm{def:g12:sound-acoustics:timbre}{espectro} é inútil para ele?
  \item Compare o seu $h$ com \qty{6.63e-34}{J.s}: qual é o erro relativo, em
        porcentagem?
\end{enumerate}

\medskip
\noindent\textbf{Parte III --- A lâmpada de hidrogênio de calibração.}
A escala de \omterm{def:g12:mechanical-waves:wavelength}{comprimentos de onda} do kit é conferida contra uma lâmpada de hidrogênio,
$E_n = -13.6/n^2$ \unit{eV}.
\begin{enumerate}[resume]
  \item Calcule $E_2$, $E_3$, $E_4$ e $E_5$.
  \item Calcule as \omterm{def:g10:energy-conservation:energy}{energias} dos \omterm{def:g12:quantum-world:photon}{fótons} das transições $3 \to 2$, $4 \to 2$ e
        $5 \to 2$.
  \item Deduza os três \omterm{def:g12:mechanical-waves:wavelength}{comprimentos de onda} e dê a cor de cada linha.
  \item Por que a lâmpada emite \emph{linhas} em vez de um arco-íris contínuo? Em uma
        frase, nomeando a ideia-chave.
  \item O hidrogênio da lâmpada conseguiria absorver luz em \qty{500}{nm}? Justifique
        com a escada.
\end{enumerate}

\medskip
\noindent\textbf{Parte IV --- Inspecionar as células: a \omterm{def:g10:signals-and-waves:wave}{onda} do elétron.}
Microfissuras numa célula ficam muito abaixo do \omterm{def:g12:projectile-motion:range}{alcance} de \qty{0.5}{\micro m} dos
microscópios ópticos; o microscópio eletrônico de varredura do laboratório opera a
\qty{10}{kV}.
\begin{enumerate}[resume]
  \item Calcule a velocidade dos elétrons; verifique que ela fica abaixo de $0.25\,c$
        (fórmula clássica).
  \item Calcule a \omterm{def:g12:newtons-laws:momentum}{quantidade de movimento} e o \omterm{def:g12:quantum-world:debroglie}{comprimento de onda de de Broglie} deles.
  \item A resolução acompanha o \omterm{def:g12:mechanical-waves:wavelength}{comprimento de onda}: calcule o ganho em relação à luz
        verde (\qty{550}{nm}); conclua com os dois números do fim de semana --- o seu $h$
        e esse ganho.
\end{enumerate}
\end{problem}
